% general commands
\clubpenalty = 10000
\widowpenalty = 10000
\displaywidowpenalty = 10000


\renewcommand*\thelstnumber{\arabic{lstnumber}.}
\lstdefinestyle{base}{
  moredelim=[is][\color{red}]{/}{/},
captionpos=b,language=C++,basicstyle=\ttfamily,
  numbers=left,numberstyle=\scriptsize\bfseries,xleftmargin=.038\textwidth
}
\lstdefinestyle{base2}{
captionpos=b,language=C++,basicstyle=\ttfamily\scriptsize,
  numbers=left,numberstyle=\scriptsize\bfseries,xleftmargin=.038\textwidth
}


\newcommand{\naive}{na\"{\i}ve }

% SETTINGS FOR EPIGRAPH (THE VIP QUOTATIONS)
\setlength\epigraphwidth{0.9\linewidth}
\setlength{\epigraphrule}{0mm} % REMOVE BAR SEPARATING EPIGRAPH FROM SOURCE
\renewcommand{\textflush}{flushepinormal} % JUSTIFY EPIGRAPH TEXT
\renewcommand{\epigraphflush}{flushright} % CENTRALIZE
\renewcommand{\sourceflush}{flushright} % CENTRALIZE

% FOR PRETTIER CANCEL
\usetikzlibrary{calc}
\tikzstyle{nosep}=[inner sep=0pt, outer sep=0pt]
\newcommand{\hcancel}[1]{%
    \tikz[baseline=(tocancel.base)]{
        \node[nosep] (tocancel) {$#1$};
        \node[nosep, yshift=.5ex]  (from) at (tocancel.south west) {};
        \node[nosep, yshift=-.5ex] (to)   at (tocancel.north east) {};
        \draw[black] (from) -- (to);
    }%
}%

%\newcommand{\bibsetup}{\thispagestyle{plain}\pagestyle{fancy}}

% UNIVERSAL PLOT SIZE
\newcommand{\plotsize}{0.47\textwidth}

% TABLE OF CONTENTS WITH EMPTY PAGE STYLE
\AtBeginDocument{\addtocontents{toc}{\protect\thispagestyle{plain}}}

% keeps the distance between paragraphs constant
\setlength{\parskip}{1ex plus 0.0ex minus 0.0ex}

\renewcommand{\floatpagefraction}{0.9}
\renewcommand{\topfraction}{0.9}
\renewcommand{\bottomfraction}{0.9}
\renewcommand{\textfraction}{0.1}
\renewcommand{\textfloatsep}{5mm}
\renewcommand{\baselinestretch}{1.15}

\newcommand*\Hide{%
\titleformat{\chapter}[display]
  {}{}{0pt}{\huge}
\titleformat{\thispagestyle{empty}\part}
  {}{}{0pt}{}
}

\newcommand\blfootnote[1]{%
  \begingroup
   \renewcommand\thefootnote{}\footnote{#1}%
   \addtocounter{footnote}{-1}%
   \endgroup
}


\newcommand{\one}{\ding{172}}
\newcommand{\two}{\ding{173}} \newcommand{\three}{\ding{174}}
\newcommand{\four}{\ding{175}} \newcommand{\five}{\ding{176}}
\newcommand{\six}{\ding{177}} \newcommand{\seven}{\ding{178}}
\newcommand{\eight}{\ding{179}} \newcommand{\nine}{\ding{180}}
\newcommand{\ten}{\ding{181}}

%\makeindex

%%%%%%%%%%%%%%%%%%%%%%%%%%%%%%%%%%%%%%%%%%%%%%%%%%%%%%%%%%%%%%%%%%%%%
% define \timestamp
\newcount\minute \newcount\hour \newcount\hourMins
\def\now{\minute=\time \hour=\time \divide \hour by 60 \hourMins=\hour \multiply\hourMins by 60
  \advance\minute by -\hourMins \zeroPadTwo{\the\hour}:\zeroPadTwo{\the\minute}}
\def\timestamp{\today\ \now}
\def\today{\the\year-\zeroPadTwo{\the\month}-\zeroPadTwo{\the\day}}
\def\zeroPadTwo#1{\ifnum #1<10 0\fi #1}
%%%%%%%%%%%%%%%%%%%%%%%%%%%%%%%%%%%%%%%%%%%%%%%%%%%%%%%%%%%%%%%%%%%%%

%\DeclareUnicodeCharacter{00A0}{ }
\DeclarePairedDelimiter{\ceil}{\lceil}{\rceil}

\newcommand{\titlepagefont}{\sffamily} % FONT USED IN THE TITLE PAGE -- change to \rmfamily FOR serif


\renewcommand*{\=}{{\kern0.1em=\kern0.1em}}
\renewcommand*{\-}{{\kern0.1em-\kern0.1em}}
\newcommand*{\+}{{\kern0.1em+\kern0.1em}}

\newenvironment{Tabular}[2][1]
  {\def\arraystretch{#1}\tabular{#2}}
  {\endtabular}
\def\changemargin#1#2{\list{}{\rightmargin#2\leftmargin#1}\item[]}
\let\endchangemargin=\endlist



%%%%%%%%%%%%%%%%%%%%%%%%%%%%%
%%%%%%%%  CODE LISTINGS %%%%%
%%%%%%%%%%%%%%%%%%%%%%%%%%%%%

\lstdefinestyle{python}{
  backgroundcolor=\color{black!3},
  belowcaptionskip=0\baselineskip,
  belowskip=0pt,
  breaklines=true,
  xleftmargin=0pt,
  language=Python,
  showstringspaces=false,
  basicstyle=\footnotesize\ttfamily,
    numbers=left,                   % where to put the line-numbers
  numberstyle=\tiny\color{Blue},  % the style that is used for the line-numbers
  stepnumber=1,                   % the step between two line-numbers. If it is 1, each line
                                  % will be numbered
  numbersep=5pt,                  % how far the line-numbers are from the code
  keywordstyle=\bfseries\color{green!40!black},
  commentstyle=\itshape\color{purple!20!black},
  identifierstyle=\color{blue!60!black},
  stringstyle=\color{orange!60!black},
  columns=flexible,
  rulecolor=\color{black},
  keepspaces=true
}
\lstset{columns=flexible}
\lstset{keepspaces=true}

\lstdefinestyle{R}{ 
  language=R,                     % the language of the code
  basicstyle=\footnotesize\ttfamily, 	  % the size of the fonts that are used for the code
  numbers=left,                   % where to put the line-numbers
  numberstyle=\tiny\color{Blue},  % the style that is used for the line-numbers
  stepnumber=1,                   % the step between two line-numbers. If it is 1, each line
                                  % will be numbered
  numbersep=5pt,                  % how far the line-numbers are from the code
  backgroundcolor=\color{black!3},  % choose the background color. You must add \usepackage{color}
  showspaces=false,               % show spaces adding particular underscores
  showstringspaces=false,         % underline spaces within strings
  showtabs=false,                 % show tabs within strings adding particular underscores
  frame=single,                   % adds a frame around the code
  rulecolor=\color{black},        % if not set, the frame-color may be changed on line-breaks within not-black text (e.g. commens (green here))
  tabsize=2,                      % sets default tabsize to 2 spaces
  captionpos=b,                   % sets the caption-position to bottom
  breaklines=true,                % sets automatic line breaking
  breakatwhitespace=false,        % sets if automatic breaks should only happen at whitespace
  keywordstyle=\color{RoyalBlue},      % keyword style
  commentstyle=\color{YellowGreen},   % comment style
  stringstyle=\color{ForestGreen}      % string literal style
}

%%%%%%%%%%%%%%%%%%%%%%%%%%%%%
%%%  TIKZ
\usetikzlibrary{calc,3d}
\usetikzlibrary{decorations.pathreplacing}


\tikzset{
    position/.style args={#1:#2 from #3}{
        at=(#3.#1), anchor=#1+180, shift=(#1:#2)
    },
    minimum size=10pt,
    every loop/.style={min distance=20mm,in=45,out=-45,looseness=10}
}

%%   COLORS AND STYLES
\definecolor{Gray}{rgb}{0.9,0.9,0.9} % THIS IS USED FOR TABLE ROWS
\definecolor{DarkGray}{rgb}{0.5,0.5,0.5} % THIS IS USED FOR PSEUDOCODE COMMENTS
\definecolor{redDark}{HTML}{9A110A}
\definecolor{redMid}{HTML}{D8140A}
\definecolor{redLight}{HTML}{FF4137}

\definecolor{blueDark}{HTML}{0B4262}
\definecolor{blueMid}{HTML}{0E5C8A}
\definecolor{blueLight}{HTML}{389CD7}

\definecolor{LightCyan}{rgb}{0.88,1,1}

\tikzstyle{article}=[draw,circle,redDark, fill=redLight, inner sep=0pt, minimum width=2mm]
\tikzstyle{author}=[draw,circle,blueDark, fill=blueLight, inner sep=0pt, minimum width=2mm]
\tikzstyle{cit}=[-latex,redMid,bend right]
\tikzstyle{collab}=[draw,blueMid]
\tikzstyle{bipart}=[draw,blueMid]

%%%%%%%%%%%%%%%%%%%%%%%%%%%%%
%%%  DATE FORMATTER        %%
%%%%%%%%%%%%%%%%%%%%%%%%%%%%%

\newcommand{\ymd}[3]{\leavevmode\hbox{#1\,/\,#2\,/\,#3}}


%%%%%%%%%%%%%%%%%%%%%%%%%%%%%
%%%  FROM SPECIAL_COMMANDS %%
%%%%%%%%%%%%%%%%%%%%%%%%%%%%%

\newenvironment{itquote}
  {\begin{quote}\itshape}
  {\end{quote}\ignorespacesafterend}


% TABLE FORMATING
\newcommand{\titlehline}{\hline\hline}


%%%%%%%%%%%%%%%%%%%%%%%%%%%%%
%%%  HYPOTHESIS %%
%%%%%%%%%%%%%%%%%%%%%%%%%%%%%
\newtheoremstyle{hypothesis}{5pt}{5pt}{\itshape}{1cm}{\bfseries}{}{.5em}{}
\theoremstyle{hypothesis}
\newtheorem{hypo}{Hypothesis}

%%%%%%%%%%%%%%%%%%%%%%%%%%%%%
%%%  MISC TYPESETTING      %%
%%%%%%%%%%%%%%%%%%%%%%%%%%%%%
\newcommand*{\refleft}{\emph{(left)}\@\xspace}
\newcommand*{\refright}{\emph{(right)}\@\xspace}


\usepackage{subcaption}
\usepackage{placeins}
\usepackage{siunitx} % Provides the S column type

\sisetup{
  round-mode = places, % Rounds numbers
  round-precision = 3, % to 3 places
  detect-all, % For aligning numbers by decimal point within table
}
